\section{Алгоритм Форда-Беллмана}

Теперь давайте решим ту же самую задачу нахождения расстояния от $s$ до всех остальных с помощью динамического программирования -- алгоритма Форда-Беллмана. В отличие от алгоритма Дейкстры этот алгоритм расширяет образ весовой функции до всей числовой прямой. Однако необходимо требовать, чтобы в графе не было циклов отрицательного веса.
Заведём двумерную динамику $d[k][v]$ -- длина кратчайшего пути из $s$ в $v$, состоящего из $k$ рёбер. Тогда легко видеть, что формула перехода представляет собой следующее:
\begin{equation*}
    d[k][v] = \min\limits_{(u, v) \in E} [ d[k - 1][u] + w(u, v) ]
\end{equation*}
То есть мы должны на каждом шаге динамики <<расширять>> путь до $v$ из смежных к ней вершин. База динамики очевидна: $dp[0][s] = 0, dp[0][v] = +\infty \, \forall v \in V \backslash \{s\}$. А ответ тогда вычисляется как: $\delta(s, t) = \min\limits_{k \in 0 \dots |V| - 1} dp[k][t]$.

\begin{algorithm}[H]
\caption{Алгоритм Форда-Беллмана}
\begin{algorithmic}[1]

\Function{Ford\_Bellman}{$G, w, s$}
    \State $dp[0][s] \gets 0$
    \State $dp[0][v] \gets +\infty\, \forall v \in V \backslash \{s\}$
    \For{$k \in 1\dots |V| - 1$}
        \For{$(u, v) \in E$}
            \State $dp[k][v] \gets \min(dp[k][v], dp[k-1][u] + w(u, v))$
        \EndFor
    \EndFor

    \State \Return $dp$
\EndFunction
\end{algorithmic}
\end{algorithm}

Время работы данного алгоритма $\Theta(|V| |E|)$.

Что если в вышеописанном алгоритме мы сделаем еще одну итерацию ($|V|$-ую)? А точнее, что если до какой-то вершины уменьшится расстояние в $dp$? Коль скоро мы рассматривали уже пути длины $|V|$, то эти пути точно самопересекаются (ведь максимальная длина несамопересекающегося пути -- это ровно $|V| - 1$). Отсюда напрашивается важное утвереждение.


\Th $\exists v \in V: \min\limits_{k \in 0 \dots |V| - 1} dp[k][v] < dp[v][|V|] \iff$ существует цикл отрицательного веса, достижимый из $s$.

\begin{proof}
    Пусть $\exists v \in V: \min\limits_{k \in 0 \dots |V| - 1} dp[v][k] < dp[v][|V|]$. Значит $\exists P = \{(s, u_1), (u_1, u_2) \dots (u_n, v)\}$ -- самопересекающийся путь $\Rightarrow$ хотя бы один цикл существует. Предпопложим, что все циклы неотрицательны.

    В силу того, что расстояние уменьшилось на $|V|$-ой итерации, то $w(P) < \min\limits_{k \in 0 \dots |V| - 1} dp[k][v]$. Далее, на основе пути $P$ построим путь $P'$. Приведём алгоритм построения $P'$, пусть $U$ -- множество раскрытых вершин (изначально пустое) и $k$ -- номер рассматриваемоего ребра (изначально $k = 0$, то есть рассматриваем $(s, u_1)$).

    \begin{enumerate}
        \item Если $k > n$, то алгоритм считаем оконченым. Рассмотрим $k$-ое ребро $(u_f, u_t)$. Если $u_t = v$, то добавляем $(u_f, u_t)$ в конец $P'$ и алгоритм завершаем. Если $u_t \in U$, то перейти к шагу 2, иначе 3.
        \item Удаляем с конца $P'$ ребра $(u_a, u_b)$, пока $u_b \neq u_t$, где $u_t$ -- вершина из пункта 1. После инкрементируем $k$ и переходим к шагу 1.
        \item Добавляем конец $P'$ ребро $(u_f, u_t)$ и в множество $U$ вершину $u_t$, инкрементируем $k$ и переходим к шагу 1.
    \end{enumerate}

    По сути алгоритм строит $P'$ и как только он видит, что путь зациклился, он удаляет цикл удалением рёбер из пути. Далее рассмотрим важные свойства полученного пути:
    \begin{itemize}
        \item $P'$ -- путь от $s$ до $v$.
        \item $P'$ -- простой. Действительно, в силу того, что мы поддерживаем множество раскрытых вершин, то все вершины в пути различны.
        \item $|P'| \leq |V| - 1$, как следствие простоты пути.
        \item $w(P') \leq w(P)$. В силу того, что мы удаляем циклы из пути $P$, которые являются неотрицательными, то легко видеть, что вес пути $P'$ уж точно не превосходит вес $P$.
    \end{itemize}
    
    Далее, $|P'| \leq |V| - 1 \Rightarrow \exists k' < |V| - 1: dp[k'][v] \leq w(P')$. Теперь заметим следующее:
    \begin{equation*}
        \min\limits_{k \in 0 \dots |V| - 1} dp[k][v] \leq dp[k'][v] \leq w(P') \leq w(P) < \min\limits_{k \in 0 \dots |V| - 1} dp[k][v]
    \end{equation*}
    \begin{equation*}
        \min\limits_{k \in 0 \dots |V| - 1} dp[k][v] < \min\limits_{k \in 0 \dots |V| - 1} dp[k][v]
    \end{equation*}
    Получили противоречие.

    В другую сторону, пусть существует цикл отрицательного веса $C = \{c_1, \dots c_k\}, k \leq |V| - 1$ достижимый из $s$. Значит на $|V|$-й итерации есть вершина цикла, которую мы рассмотрим второй раз вдоль цикла, а значит мы обновим минимум.
\end{proof}

\section{Алгоритм Флойда-Уоршелла}

Всюду далее мы опять считаем, что образ весовой функции вещественен (не только положительный). Теперь наша задача состоит в том, чтобы найти расстояние между всеми парами вершин.

Давайте пронумеруем все вершины $V = \{v_1, \dots, v_{|V|}\}$ и заведём динамику $dp[k][u][v]$ -- минимальный вес пути из $u$ в $v$ по \textit{промежуточным} вершинам с номерами не больших $k$.
Далее воспользуемся следующей идеей:
\begin{equation*}
    \delta(u, v) = \min\limits_{p \in V} [\delta(u, p) + \delta(p, v)]
\end{equation*}
Будем делать переход $dp[k][u][v] = \min(dp[k - 1][u][v_k] + dp[k - 1][v_k][v], dp[k - 1][u][v])$, по сути пытаясь выбрать оптимальную <<среднюю вершину>>.

\begin{algorithm}[H]
\caption{Алгоритм Флойда-Уоршелла}
\begin{algorithmic}[1]

\Function{Floyd\_Warshall}{$G, w$}
    \State $dp[0][u][u] \gets 0, \forall u \in V$
    \State $dp[0][u][v] \gets w(u, v), \forall u, v \in V$
    \State $dp[k][u][v] \gets +\infty, \forall k \in [1, n], \forall u \neq v \in V$
    \For{$k \in 1\dots |V|$}
        \For{$u \in V$}
            \For{$v \in V$}
                \State $dp[k][u][v] \gets \min(dp[k - 1][u][v_k] + dp[k - 1][v_k][v], dp[k - 1][u][v])$
            \EndFor
        \EndFor
    \EndFor
    \State \Return $dp$
\EndFunction
\end{algorithmic}
\end{algorithm}
Очевидно, ответ будет в $dp[|V|][u][v]$. Асимптотика $\Theta\left(|V|^3\right)$.

Аналогично алгоритму Форда-Беллмана здесь тоже можно искать циклы отрицательного веса.

\Th В графе есть циклы отрицательного веса $\iff \exists v \in V: dp[|V|][v][v] < 0$
\begin{proof}
    $\exists v \in V: dp[|V|][v][v] < 0 \iff \exists k \in 1 \dots |V|: dp[|V|][v][v] = dp[|V| - 1][v][v_k] + dp[|V| - 1][v_k][v] < 0 \iff$ В графе есть цикл отрицательного веса.

    \textit{Пояснение:} во второй эквивалентности мы по сути утверждаем, что существует такая вершина $v_k$, что путь из $v$ в $v$ (то есть цикл) будем отрицательного веса, поэтому третья эквивалентность верна.
\end{proof}
